\documentclass[11pt]{article}

\usepackage[T1]{fontenc}
\usepackage[utf8]{inputenc}
\usepackage{amsmath,amssymb,amsthm}
\usepackage{newtxtext,newtxmath}
\usepackage[margin=1.15in]{geometry}
\usepackage{booktabs}
\usepackage{microtype}
\usepackage{xcolor}
\usepackage{enumitem}
\definecolor{link}{HTML}{2F4A8C}
\usepackage[colorlinks=true,linkcolor=link,citecolor=link,urlcolor=link]{hyperref}

\setlist[itemize]{itemsep=2pt,topsep=4pt,parsep=0pt}
\newtheorem{proposition}{Proposition}
\newcommand{\yv}{\mathbf{y}}
\newcommand{\Cm}{\mathbf{C}}
\newcommand{\Dm}{\mathbf{D}}
\newcommand{\Idm}{\mathbf{I}}
\newcommand{\Km}{\mathbf{K}}
\newcommand{\Wm}{\mathbf{W}}
\newcommand{\one}{\mathbf{1}}
\newcommand{\thetav}{\boldsymbol{\theta}}
\newcommand{\Sig}{\boldsymbol{\Sigma}}
\newcommand{\BM}{\mathrm{BM1}}
\newcommand{\OUo}{\mathrm{OU1}}
\newcommand{\OUx}{\mathrm{OUx}}
\newcommand{\E}{\mathbb{E}}

\title{\vspace{-2em}A Bayesian formulation of SCOUT's per-gene test\\for adaptive expression evolution}
\author{\normalsize\texttt{sparseou} technical note --- specification for \texttt{stan/scout.stan}}
\date{\normalsize 25 August 2026}

\begin{document}
\maketitle
\vspace{-1.5em}

\begin{abstract}
\noindent
SCOUT fits three hypotheses to each gene on a single-cell lineage tree --- neutral drift, stabilising
selection on one optimum, and adaptation to several regime-specific optima --- and reports the
minimum-AICc winner as a single word. At small tree sizes that word is correct only slightly more
often than chance, and it carries no statement of confidence. This note specifies the same three
hypotheses as a Bayesian model comparison. The discrete model indicator is marginalised analytically,
so a single fit per gene returns posterior model probabilities together with posteriors for the
selection strength, the drift rate and the optima. The formulation is deliberately \emph{not}
hierarchical: genes are treated independently, exactly as SCOUT treats them, so that it can serve as
an honest baseline against which pooling across genes is later measured.
\end{abstract}

\section{What changes, and what does not}

Given a rooted lineage tree, a gene's expression at the tips, and a \emph{supplied} partition of the
tips into regimes, SCOUT fits three models by maximum likelihood, penalises each by its parameter
count through AICc, and returns the argmax. Two properties of that procedure motivate what follows.

First, the output is a hard label. A gene whose evidence is evenly split between two hypotheses
receives the same kind of answer as one whose evidence is overwhelming. Second, AICc is an asymptotic
criterion, and it is being applied at tree sizes as small as $n = 32$.

Replacing the argmax with a posterior addresses both: integrating over parameters rather than
maximising, and reporting $\Pr(M \mid \yv)$ rather than $\arg\max$. What does \emph{not} change is the
scientific question, the likelihood, or the requirement that the adaptive landscape be supplied.

\section{Setup and notation}

Let the tree be rooted and \textbf{ultrametric}, with $n$ tips and height $H$. Fix one gene and let
$\yv \in \mathbb{R}^{n}$ be its expression across the tips, standardised to zero mean and unit
variance. Write

\begin{center}
\begin{tabular}{@{}ll@{}}
\toprule
Symbol & Meaning \\
\midrule
$K$ & number of supplied regimes, $k = 1,\dots,K$ \\
$r(u) \in \{1,\dots,K\}$ & regime painted on the branch at time $u$ --- data, not inferred \\
$H$ & root-to-tip time, identical for every tip \\
$D_{ij} = 2\bigl(H - t_{\mathrm{mrca}(i,j)}\bigr)$ & patristic distance between tips $i$ and $j$ \\
$C_{ij} = t_{\mathrm{mrca}(i,j)}$ & shared root-to-MRCA path length \\
\bottomrule
\end{tabular}
\end{center}

\noindent
Throughout, $\Km(\alpha)$ denotes the $n \times n$ matrix with entries
\begin{equation}
\bigl[\Km(\alpha)\bigr]_{ij} \;=\; e^{-\alpha D_{ij}},
\end{equation}
which is an elementwise exponential of $-\alpha\Dm$, not a matrix exponential. Both Ornstein--Uhlenbeck
hypotheses share the covariance
\begin{equation}
\Sig(\alpha, s, \tau) \;=\; s^{2}\,\Km(\alpha) \;+\; \tau^{2}\,\Idm .
\label{eq:sig}
\end{equation}

\section{The three hypotheses}

A single discrete parameter carries the scientific question,
\begin{equation}
M \in \{\BM,\ \OUo,\ \OUx\}, \qquad \Pr(M = m) = \pi_m, \qquad \pi = (\tfrac13,\tfrac13,\tfrac13),
\end{equation}
and everything else in the model exists to give it something to weigh. All three hypotheses are
multivariate normal over the tips; they differ only in the \emph{mean} --- whether there is an
optimum at all, and whether it is the same everywhere on the tree.

\paragraph{Neutral drift.}
Brownian motion started at $\theta_{\mathrm{bm}}$, with no mean reversion, so variance accumulates
with shared ancestry:
\begin{equation}
\yv \mid M = \BM \;\sim\; \mathcal{N}_{n}\!\left(\theta_{\mathrm{bm}}\one_{n},\;
   \sigma_{\mathrm{bm}}^{2}\,\Cm + \tau_{1}^{2}\,\Idm\right).
\end{equation}

\paragraph{One global optimum.}
An Ornstein--Uhlenbeck process
\begin{equation}
dX(t) \;=\; \alpha\bigl(\theta - X(t)\bigr)\,dt \;+\; \sigma\,dW(t),
\label{eq:ou}
\end{equation}
with the root drawn from its stationary distribution about that optimum. Being stationary, the mean
is $\theta$ everywhere and the covariance depends on the tips only through patristic distance:
\begin{equation}
\yv \mid M = \OUo \;\sim\; \mathcal{N}_{n}\!\left(\theta_{1}\one_{n},\; \Sig(\alpha_{1}, s_{1}, \tau_{2})\right).
\end{equation}
The scale is parameterised by the \emph{stationary} standard deviation $s^{2} = \sigma^{2}/(2\alpha)$
rather than by $\sigma$, since $s$ is comparable across values of $\alpha$ and takes a prior on the
same footing as the standardised data; $\sigma = s\sqrt{2\alpha}$ is recovered afterwards.

\paragraph{One optimum per regime.}
The same process, but with an optimum that is piecewise constant along the tree,
$\theta(u) = \theta_{r(u)}$:
\begin{equation}
\yv \mid M = \OUx \;\sim\; \mathcal{N}_{n}\!\left(\Wm(\alpha_{3})\,\thetav,\; \Sig(\alpha_{3}, s_{3}, \tau_{3})\right),
\qquad \thetav \in \mathbb{R}^{K}.
\label{eq:oux}
\end{equation}

\paragraph{Nesting.}
$\BM$ is the $\alpha \to 0$ limit of $\OUo$, and $\OUo$ is $\OUx$ with all $\theta_k$ equal. They are
nonetheless fitted as three hypotheses with their own parameters rather than as restrictions of one,
so that each carries an uncontaminated marginal likelihood.

\section{The exposure matrix}

Integrating \eqref{eq:ou} along the root-to-tip path of tip $i$ gives
\begin{equation}
X(H) \;=\; e^{-\alpha H} X(0) \;+\; \int_{0}^{H}\!\alpha\, e^{-\alpha(H-u)}\,\theta(u)\,du
   \;+\; \sigma\!\int_{0}^{H}\! e^{-\alpha(H-u)}\,dW(u).
\end{equation}
Let $\mathcal{P}(i)$ be the decomposition of that path into maximal segments $(k, a, b)$, meaning
regime $k$ occupies the time interval $[a,b]$. With the root started at its own optimum, the
stochastic term has mean zero and
\begin{equation}
\E[y_i] \;=\; e^{-\alpha H}\theta_{r(0)}
   \;+\!\! \sum_{(k,a,b) \in \mathcal{P}(i)}\!\! \theta_{k}\left(e^{-\alpha(H-b)} - e^{-\alpha(H-a)}\right)
   \;=\; \bigl[\Wm(\alpha)\thetav\bigr]_{i}.
\label{eq:W}
\end{equation}

Each row of $\Wm$ records how much of that tip's ancestry sits in each regime, discounted by how
quickly the process forgets: a tip that entered a regime long ago has arrived at its optimum, while
one that entered recently still carries its ancestor's.

\begin{proposition}[$\Wm$ is row-stochastic]
For every tip $i$ and every $\alpha > 0$, $\sum_{k} W_{ik}(\alpha) = 1$.
\end{proposition}
\begin{proof}
The segments of $\mathcal{P}(i)$ tile $[0,H]$ contiguously, so the sum in \eqref{eq:W} telescopes to
$e^{-\alpha(H-H)} - e^{-\alpha(H-0)} = 1 - e^{-\alpha H}$. Adding the root term $e^{-\alpha H}$ gives
unity.
\end{proof}

\noindent
The tip mean is therefore a genuine weighted average of the optima. This is worth verifying
numerically rather than assuming: it fails immediately if the regime painting or the segment
decomposition is wrong, and it is the cheapest available check on both.

\section{Priors}

\begin{center}
\begin{tabular}{@{}lll@{}}
\toprule
Parameter & Prior & Role \\
\midrule
$\theta_{\mathrm{bm}},\,\theta_{1},\,\theta_{k}$ & $\mathcal{N}(0, s_{\theta}^{2})$ & optima --- see \S\ref{sec:occam} \\
$\sigma_{\mathrm{bm}},\, s_{1},\, s_{3}$ & $\mathcal{N}^{+}(0, s_{\mathrm{scale}}^{2})$ & drift rate / stationary spread \\
$h_{1},\, h_{3}$ & $\mathrm{LogNormal}(0, s_{h}^{2})$ & phylogenetic half-life, in tree heights \\
$\tau_{1},\, \tau_{2},\, \tau_{3}$ & $\mathcal{N}^{+}(0, s_{\tau}^{2})$ & measurement error \\
\bottomrule
\end{tabular}
\end{center}

\noindent
Selection strength is parameterised by half-life relative to tree height rather than by $\alpha$
directly,
\begin{equation}
\alpha_{m} \;=\; \frac{\ln 2}{h_{m} H}, \qquad m \in \{\OUo, \OUx\},
\end{equation}
so that $h = 1$ reads as ``the optimum is reached over the tree's lifetime''. The reparameterisation
keeps the sampler away from both degenerate ends: $\alpha \to 0$, where the process is
indistinguishable from drift, and $\alpha \to \infty$, where it is indistinguishable from white noise.

\section{Inference}

Stan cannot sample a discrete parameter, so the indicator is summed out. Writing $\phi$ for the
collection of continuous parameters and $\ell_{m}(\phi) = \log p(\yv \mid \phi, M = m)$, the sampled
log-density is
\begin{equation}
\log p(\yv, \phi) \;=\; \log \sum_{m} \exp\bigl\{\log \pi_{m} + \ell_{m}(\phi)\bigr\} \;+\; \log p(\phi).
\end{equation}

Because $\Pr(M = m \mid \yv) = \int \Pr(M = m \mid \phi, \yv)\, p(\phi \mid \yv)\, d\phi$, the model
probability is the posterior mean of a quantity available at every draw. No bridge sampling and no
separate marginal-likelihood machinery is required:
\begin{equation}
\Pr(M = m \mid \yv) \;=\; \E_{\phi \mid \yv}
   \left[\frac{\pi_{m}\exp \ell_{m}(\phi)}{\sum_{j} \pi_{j}\exp \ell_{j}(\phi)}\right].
\label{eq:pm}
\end{equation}
This is a Rao--Blackwellised estimator: the conditional probability is evaluated exactly at each
draw and averaged, rather than the indicator being sampled and its visits counted.

The output per gene is three numbers summing to one. A gene whose evidence is genuinely ambiguous
returns something near $(0.3, 0.4, 0.3)$ instead of a confident single word --- which is the entire
point of the exercise.

\section{The prior on the optima is the Occam penalty}
\label{sec:occam}

Of all the priors in the model, $s_{\theta}$ is the one that is not a nuisance. Marginalising the
optima out of \eqref{eq:oux} analytically, which is possible in closed form because the mean is
linear in $\thetav$, gives
\begin{equation}
p(\yv \mid M = \OUx, \alpha, s, \tau) \;=\; \mathcal{N}_{n}\!\left(\yv \;\middle|\; \mathbf{0},\;
   \Sig(\alpha, s, \tau) + s_{\theta}^{2}\,\Wm(\alpha)\Wm(\alpha)^{\!\top}\right).
\label{eq:occam}
\end{equation}

The penalty is now explicit. Freeing $K$ optima inflates the covariance by $s_{\theta}^{2}\Wm\Wm^{\top}$,
a rank-$K$ term, and the resulting log-determinant grows like $K \log s_{\theta}$. That is what
replaces AICc's $2k$: a larger $s_{\theta}$ improves the fit but is charged for the additional freedom.
Set $s_{\theta}$ wide and $\OUx$ can never win; set it narrow and it wins too easily.

Nothing else in the model carries comparable leverage over the answer. It is therefore exposed as
data rather than fixed internally, and a sweep across it belongs in the results rather than in a
robustness appendix. Note that \eqref{eq:occam} also suggests a more efficient implementation than
the one specified here, in which $\thetav$ is never sampled; that is deliberately deferred, since
posteriors for the individual optima are wanted as an output.

\section{Numerical considerations}

\begin{itemize}
  \item A fixed $10^{-8}\Idm$ is added to every covariance. As $\alpha \to 0$, $\Km(\alpha) \to \one\one^{\top}$,
    which has rank one; without the jitter every Cholesky factorisation fails at the low-$\alpha$ end
    of the prior.
  \item $h_{m}$ is bounded to $[10^{-4}, 10^{4}]$, keeping $\Km(\alpha)$ clear of overflow and underflow.
  \item $\BM$ is evaluated in the eigenbasis of $\Cm$. Since $\Cm$ and $\yv$ are both data, the
    decomposition and the rotation of $\yv$ are performed once, leaving only scalars to vary --- so
    $\BM$ costs $O(n)$ per gradient evaluation rather than $O(n^{3})$. The two OU hypotheses retain
    independent $\alpha$, as OUwie does, and so require a Cholesky factorisation each.
  \item Chains are started from \emph{jittered} initial values. Identical starts leave $\hat{R}$
    unable to detect between-chain disagreement: a 256-tip fit reported $\hat{R} = 1.09$ from
    identical inits and $\hat{R} = 1.01$ once jittered.
\end{itemize}

\section{Scope, and deviations from SCOUT}

Deliberately excluded: any sharing of information across genes --- that is the hierarchical model
this specification exists to be a baseline for. Regimes are supplied, never inferred. Expression is
Gaussian on log-normalised counts, with no count model. Trees are ultrametric, and the dense Cholesky
caps practical tree size near $n = 256$.

Three deviations from SCOUT are genuine and should be treated as open until checked against the
original OUwie implementation.

\begin{itemize}
  \item \textbf{The measurement-error term $\tau$.} SCOUT's smoothed arm has no such parameter,
    though their EM and tip-fog arms do. It is included here because lineage smoothing multiplies the
    data by a similarity matrix $\mathbf{S}$, giving a covariance $\mathbf{S}\Sig\mathbf{S}^{\top}$
    that no Ornstein--Uhlenbeck process can produce. A generative model fitted to smoothed data
    absorbs that mismatch as spurious signal.
  \item \textbf{Regime painting.} SCOUT paints internal branches using \texttt{ape::ace} under an
    equal-rates Mk model. The present implementation uses a deterministic parsimony rule, which
    coincides on simulated trees where clades are regime-coherent, and may not on real data.
  \item \textbf{Root treatment.} The root is placed at its own optimum here. OUwie exposes a
    \texttt{root.station} option; if their setting differs, the $\OUo$ and $\OUx$ likelihoods differ
    in a way that no amount of correctness in the Stan program would reveal.
\end{itemize}

\section{Recovery study I: neutral drift}
\label{sec:recovery-bm}

Before any model comparison is attempted, each hypothesis should be checked in isolation: simulate a
gene from it, fit that hypothesis alone, and ask whether the generating parameters come back. This
section does so for $\BM$. The gene is simulated from exactly the model being fitted, $\BM$ is fitted
alone (no competing hypotheses, no model indicator), and the data are deliberately \emph{not}
standardised, so $\sigma$, $\tau$ and $\theta$ retain their generating units and recovery is a real
check rather than a check up to an unknown rescaling.

\subsection*{What is being asked}

The generating process is Brownian motion on a height-one ultrametric tree, observed with additive
noise:
\begin{equation}
y_i \;=\; X(t_i) \;+\; \varepsilon_i, \qquad
X \sim \mathrm{BM}\bigl(\theta_{\mathrm{bm}},\, \sigma_{\mathrm{bm}}^{2}\bigr), \qquad
\varepsilon_i \overset{\text{iid}}{\sim} \mathcal{N}(0, \tau^{2}),
\end{equation}
with $\theta_{\mathrm{bm}} = 2$, $\sigma_{\mathrm{bm}} = 1$, $\tau = 0.3$. The three parameters play
quite different roles, and it is worth being explicit about what success would even look like for
each.

\paragraph{$\sigma_{\mathrm{bm}}$, the drift rate.} Estimated from the $\binom{n}{2}$ pairwise
contrasts, whose expected squared differences grow with patristic distance. Information accumulates
with $n$; this should recover sharply.

\paragraph{$\tau$, the measurement error.} Identified by \emph{sister pairs}. Two tips separated by a
very short path should be nearly identical under pure drift; the extent to which they are not is
$\tau$. A tree with many recent splits is informative about $\tau$; a star tree would not be.

\paragraph{$\theta_{\mathrm{bm}}$, the root state.} This is the interesting one, and the reason the
study is worth running. Every tip descends from one root, so the $n$ tips do not furnish $n$
independent observations of $\theta$. The generalised-least-squares variance is
\begin{equation}
\operatorname{Var}(\hat\theta_{\mathrm{bm}}) \;=\; \bigl(\one^{\!\top} V^{-1} \one\bigr)^{-1},
\qquad V \;=\; \sigma_{\mathrm{bm}}^{2}\,\Cm + \tau^{2}\,\Idm,
\label{eq:gls}
\end{equation}
which is computable before any sampling and is far larger than the na\"ive $\sqrt{\mathrm{var}/n}$.
Writing $n_{\mathrm{eff}} = (\sigma^{2} + \tau^{2}) / \operatorname{Var}(\hat\theta)$ for the number
of independent observations the tree is effectively worth:

\begin{center}
\begin{tabular}{@{}rccc@{}}
\toprule
$n$ tips & $\operatorname{sd}(\hat\theta_{\mathrm{bm}})$ from \eqref{eq:gls} & if tips were iid & $n_{\mathrm{eff}}$ \\
\midrule
 32 & 0.443 & 0.185 & 5.6 \\
 64 & 0.407 & 0.131 & 6.6 \\
128 & 0.377 & 0.092 & 7.7 \\
256 & 0.353 & 0.065 & 8.7 \\
512 & 0.333 & 0.046 & 9.8 \\
\bottomrule
\end{tabular}
\end{center}

\noindent
Sixteen-fold more data buys roughly four more effective observations: $n_{\mathrm{eff}}$ grows like
$\log_2 n$, not like $n$. This is not a defect of the estimator but a property of the design ---
a phylogeny is a deeply nested sampling scheme, and the root is the one quantity all of it shares.
It gives a sharp, pre-registered prediction: \textbf{a correct implementation must return a posterior
for $\theta_{\mathrm{bm}}$ roughly four times wider than iid intuition suggests, and that width must
barely shrink as $n$ grows.} An implementation that reported $\operatorname{sd} \approx 0.09$ at
$n = 128$ would be confidently wrong, and no amount of posterior-predictive checking on the tips
would reveal it.

\subsection*{What Stan returns}

$\BM$ was fitted alone to one simulated dataset at each of four tree sizes, four chains, 1000
warm-up and 1000 sampling iterations. Every fit reported $\hat{R} \le 1.005$ and zero divergent
transitions, and each took under $1.3$ seconds.

\begin{center}
\begin{tabular}{@{}rccccc@{}}
\toprule
$n$ & $\theta_{\mathrm{bm}}$ (truth $2.0$) & $\sigma_{\mathrm{bm}}$ (truth $1.0$) & $\tau$ (truth $0.3$)
  & $\operatorname{sd}(\theta)$ predicted & $\hat{R}$ \\
\midrule
 32 & $1.89 \pm 0.28$ & $0.55 \pm 0.27$ & $0.50 \pm 0.16$ & $0.443$ & $1.003$ \\
 64 & $2.65 \pm 0.45$ & $1.09 \pm 0.16$ & $0.27 \pm 0.15$ & $0.407$ & $1.002$ \\
128 & $2.41 \pm 0.376$ & $1.008 \pm 0.114$ & $0.23 \pm 0.11$ & $\mathbf{0.377}$ & $1.005$ \\
256 & $1.77 \pm 0.40$ & $1.15 \pm 0.08$ & $0.16 \pm 0.09$ & $0.353$ & $1.004$ \\
\bottomrule
\end{tabular}
\end{center}

Two things in that table matter more than the fact that the intervals contain the truth.

First, at $n = 128$ the posterior standard deviation of $\theta_{\mathrm{bm}}$ is $0.3756$ against the
$0.377$ predicted by \eqref{eq:gls} --- agreement to $0.4\%$ on a quantity computed before sampling
began. That is a far stronger check than coverage of a point value, because it tests the covariance
structure rather than just the location.

Second, the $\operatorname{sd}(\theta)$ column is effectively flat across an eightfold increase in
tips while the iid expectation falls from $0.185$ to $0.065$. The predicted non-shrinkage is visible
in the fitted output, not merely in the algebra.

At $n = 32$ the fit is genuinely poor: $\sigma = 0.55$ against a true $1.0$, with $\tau$
correspondingly inflated to $0.50$. This is the drift--noise trade-off, not a bug. Both parameters
add variance, and $32$ tips do not contain enough short branches to separate a slow diffusion from a
noisy observation of a fast one. Across $150$ replicate datasets the posterior means of
$\sigma_{\mathrm{bm}}$ and $\tau$ correlate at $r = -0.73$.

\subsection*{Calibration}

Recovery on a single dataset cannot distinguish a correct implementation from a lucky one. The model
was therefore refitted to $150$ independently simulated datasets at $n = 128$, asking how often the
central credible intervals contain the generating value. Nominal coverage should be recovered within
Monte Carlo error, and is:

\begin{center}
\begin{tabular}{@{}lccc@{}}
\toprule
Parameter & $50\%$ interval & $90\%$ interval & mean posterior mean \\
\midrule
$\theta_{\mathrm{bm}}$ & $0.467 \pm 0.080$ & $0.853 \pm 0.057$ & $1.935$ \quad (truth $2.0$) \\
$\sigma_{\mathrm{bm}}$ & $0.520 \pm 0.080$ & $0.927 \pm 0.042$ & $1.007$ \quad (truth $1.0$) \\
$\tau$                 & $0.520 \pm 0.080$ & $0.873 \pm 0.053$ & $0.265$ \quad (truth $0.3$) \\
\bottomrule
\end{tabular}
\end{center}

\noindent
All six figures sit within two standard errors of nominal, and no fit in the $150$ produced a single
divergent transition. The one visible bias is in $\tau$, whose posterior mean averages $0.265$ against
a true $0.3$: the half-normal prior pulls a small positive parameter toward zero. Interval coverage is
unaffected, so this is prior shrinkage behaving as specified rather than an error.

\paragraph{Scope of this check.} Coverage is evaluated at one fixed parameter value, not averaged over
the prior, so it is a frequentist calibration statement rather than simulation-based calibration.
Full SBC --- drawing the generating parameters from the prior at each replicate --- is the stronger
test and is now cheap, since a fit takes well under a second.

\section{Validation status}

\begin{center}
\begin{tabular}{@{}llp{6.6cm}@{}}
\toprule
Check & Status & Result \\
\midrule
$\Wm(\alpha)$ vs.\ reference construction & done & agrees to $10^{-16}$; rows sum to $1$ \\
$\BM$ likelihood vs.\ independent ML & done & exact to four decimals \\
$\OUo$, $\OUx$ vs.\ independent ML & done & agree to $0.005$ where the optimiser converges \\
Sampling behaviour at $n = 256$ & done & $\hat{R} \le 1.02$, ESS $\ge 250$, jittered inits \\
Independent ML vs.\ original OUwie & \emph{open} & gates everything below \\
Accuracy and calibration on truth-labelled data & \emph{open} & \\
Behaviour across tree size & \emph{open} & \\
\bottomrule
\end{tabular}
\end{center}

\noindent
The first four checks establish that the program computes the intended likelihood. They say nothing
about whether that likelihood is SCOUT's, which is the purpose of the fifth --- and which matters,
because a faithful reimplementation of a drifted reimplementation is still drifted.

\end{document}
